\documentclass[runningheads,a4paper,14pts]{llncs}
\usepackage[T1]{fontenc}
\usepackage[utf8]{inputenc}
\usepackage[spanish,es-noshorthands]{babel} 
\usepackage{amsmath,amssymb}
\usepackage{graphicx}
\usepackage{booktabs}
\usepackage{array}
\usepackage{xcolor}
\usepackage{url}

\usepackage{tikz}
\usetikzlibrary{positioning,arrows.meta,shapes.geometric,fit,backgrounds,calc}
\tikzset{
	stage/.style={draw, rounded corners=2pt, align=center, font=\scriptsize,
		fill=blue!6, minimum height=7mm, inner sep=3pt},
	store/.style={draw, align=center, font=\scriptsize, fill=orange!12,
		cylinder, shape border rotate=90, aspect=0.25,
		minimum height=9mm, inner sep=2pt},
	hook/.style={draw, dashed, rounded corners=2pt, align=center, font=\scriptsize\itshape,
		fill=green!8, minimum height=6mm, inner sep=2pt},
	flow/.style={-{Stealth[length=2mm]}, thick},
}

\newcommand{\keywords}[1]{\par\addvspace\baselineskip
	\noindent\keywordname\enspace\ignorespaces#1}
\addto\captionsspanish{\renewcommand{\abstractname}{Abstract}}
\graphicspath{{figuras/}}

\begin{document}
	
	\mainmatter
	
	\title{ShealtRI--YourHealthWiki: un Sistema de Recuperación de Información Médica basado en Indexado Semántico Latente y Generación Aumentada por Recuperación}
	
	\titlerunning{ShealtRI: Recuperación de Información Médica con LSI y RAG}
	
	\author{Lianny de la C. Reveé Valdivieso\inst{1} \and Kevin A. Torres Perera\inst{1} \and Jocdan L. López Mantecón\inst{1}}
	
	\authorrunning{L. Reveé, K. Torres, J. López}
	
	\institute{Ciencias de la Computación, Grupos 311 y 312, Universidad de La Habana\\
		\url{https://github.com/MatCom-Developing-Team-10/ShealtRI-YourHealthWiki.git}}
	
	\toctitle{ShealtRI--YourHealthWiki}
	\maketitle
	
	\begin{abstract}
		This paper presents \emph{ShealtRI--YourHealthWiki}, an information retrieval (IR)
		system for the health and medicine domain that lets users query medical
		information in natural language and returns relevant documents together with a
		profile-aware, retrieval-augmented answer. The non-basic retrieval model is
		Latent Semantic Indexing (LSI): a TF--IDF term--document matrix is reduced with
		truncated singular value decomposition (SVD) to a latent space of $k=100$
		dimensions, and the resulting latent vectors are used as embeddings inside a
		ChromaDB vector store searched by cosine similarity. The architecture is a
		Pipeline+Microkernel modular monolith: a fixed retrieval pipeline (parsing,
		spell correction, retrieval, hybrid re-ranking, generation) is extended by
		optional plugins (thesaurus-based query expansion, Rocchio relevance feedback,
		a content-based recommender, and an evaluation module) that attach through
		declarative hooks without contaminating the mandatory code. A sitemap-driven
		crawler acquires curated content from MedlinePlus, Mayo Clinic and the NHS while
		honouring \texttt{robots.txt}, per-domain rate limiting and retry policies. When
		the local corpus is insufficient---detected through result-count, score and
		recency criteria---an automatic DuckDuckGo web fallback supplements and is cached
		for reuse. Retrieved documents are re-ranked by a hybrid LSI+BM25 scorer and
		passed to a Groq-hosted large language model grounded strictly on the retrieved
		context to mitigate hallucination, with a template fallback when no model is
		available. We describe the end-to-end data flow, the two-level storage design,
		the dynamic-indexing (folding-in) mechanism, and an evaluation module reporting
		Precision, Recall, F1, NDCG, MAP and MRR. We close with the detected limitations
		and the design we would pursue under idealised conditions.
		\keywords{recuperación de información, indexado semántico latente, LSI, TF--IDF,
			RAG, retroalimentación de Rocchio, dominio médico}
	\end{abstract}
	

	\setcounter{tocdepth}{4}
	\tableofcontents
	\newpage
	
	%==================================================================
	\section{Introducción}
	%==================================================================
	
	Los Sistemas de Recuperación de Información (SRI) constituyen la disciplina que
	estudia la representación, el almacenamiento, la organización y el acceso a
	ítems de información para satisfacer una necesidad expresada por un usuario
	\cite{manning2008,baeza1999}. Su relevancia es hoy mayor que nunca: el volumen
	de texto disponible crece a un ritmo que vuelve impracticable la búsqueda
	puramente manual, y la calidad de las decisiones ---especialmente en dominios
	sensibles como la salud--- depende de poder localizar, en segundos, contenido
	fiable y pertinente. A diferencia de un sistema de bases de datos, que devuelve
	coincidencias exactas sobre datos estructurados, un SRI opera sobre lenguaje
	natural intrínsecamente ambiguo y debe ordenar los resultados por relevancia
	estimada, no por igualdad.
	
	Este informe documenta \emph{ShealtRI--YourHealthWiki}, un SRI para el dominio de
	\textbf{salud y medicina} desarrollado como proyecto integrador. El sistema
	permite formular consultas en lenguaje natural sobre síntomas, enfermedades,
	medicamentos y cuidados, y recupera documentos médicos relevantes empleando
	\emph{Latent Semantic Indexing} (LSI) como modelo de recuperación no básico,
	enriqueciendo la respuesta mediante \emph{Retrieval-Augmented Generation} (RAG)
	adaptada al perfil del usuario. Por un lado, ofrecer una
	experiencia de búsqueda médica que tolere errores ortográficos y sinonimia ---
	frecuentes cuando un paciente describe su problema con vocabulario no técnico---;
	por otro, garantizar que la respuesta generada esté \emph{anclada} en fuentes
	recuperadas verificables, evitando la fabricación de información.
	
	El alcance abarca todo el ciclo de un SRI: adquisición de datos mediante un
	\emph{crawler} respetuoso de las políticas de los sitios, indexación con índice
	invertido y reducción semántica, recuperación y re-ordenamiento híbrido,
	generación de respuestas, y un módulo de evaluación cuantitativa.
	
	%==================================================================
	\section{Marco Teórico}\label{sec:teoria}
	%==================================================================
	
	\subsection{El modelo formal de un SRI}
	
	Siguiendo la formalización clásica \cite{salton1983,baeza1999}, un SRI se define
	como una cuádrupla $\langle D, Q, F, R\rangle$. En ShealtRI, $D$ es la colección
	de documentos médicos representados como vectores en un espacio semántico latente $\mathbb{R}^{k}$; $Q$
	es el conjunto de consultas en lenguaje natural, normalizadas, corregidas
	ortográficamente y proyectadas al mismo espacio; $F$ es el \emph{framework} de
	emparejamiento, esto es, el modelo de espacio vectorial latente con similitud del
	coseno, complementado por una fusión léxica con BM25; y $R$ es la función de
	ranking $R(q,d)\!\in\![0,1]$ que ordena los documentos por relevancia estimada
	para cada consulta.
	
	\subsection{Modelo de espacio vectorial y ponderación TF--IDF}
	
	El modelo de espacio vectorial representa documentos y consultas como vectores
	sobre el vocabulario del corpus \cite{salton1983}. La componente asociada a un
	término $t$ en un documento $d$ se pondera con el esquema TF--IDF,
	\begin{equation}
		w_{t,d} = \mathrm{tf}_{t,d}\cdot \log\frac{N}{\mathrm{df}_t},
	\end{equation}
	donde $\mathrm{tf}_{t,d}$ es la frecuencia del término en el documento, $N$ el
	número de documentos y $\mathrm{df}_t$ la frecuencia documental de $t$. TF--IDF
	premia los términos frecuentes en un documento pero raros en la colección, que
	son los más discriminantes \cite{manning2008}. La relevancia entre consulta $q$
	y documento $d$ se mide con la similitud del coseno,
	$\mathrm{sim}(q,d)=\frac{\mathbf{q}\cdot\mathbf{d}}{\lVert\mathbf{q}\rVert\,\lVert\mathbf{d}\rVert}$.
	
	\subsection{Indexado Semántico Latente (LSI): el modelo no básico}
	
	El modelo de espacio vectorial puro sufre dos problemas léxicos bien conocidos:
	la \emph{sinonimia} (términos distintos con igual significado, p.\,ej.\ ``infarto''
	y ``ataque al corazón'') y la \emph{polisemia}. LSI \cite{deerwester1990} los
	mitiga aplicando una descomposición en valores singulares truncada (truncated
	SVD) sobre la matriz término--documento $A$ (de TF--IDF):
	\begin{equation}
		A \approx U_k\,\Sigma_k\,V_k^{\top},
	\end{equation}
	donde $U_k$, $\Sigma_k$ y $V_k^{\top}$ retienen las $k$ mayores componentes. Las
	columnas de $V_k^{\top}$ son las representaciones de los documentos en un espacio
	\emph{latente} de $k$ dimensiones que captura relaciones de co-ocurrencia entre
	términos, de modo que documentos semánticamente afines quedan cercanos aunque no
	compartan vocabulario exacto. Una consulta se proyecta al mismo espacio mediante
	$\mathbf{q}_{\text{proj}} = \mathbf{q}^{\top} U_k \Sigma_k^{-1}$, y luego se
	comparan los vectores por coseno. La ventaja frente a los modelos básicos
	(booleano y vectorial puro) para el dominio médico es directa: mejora el
	\emph{recall} ante la sinonimia abundante de la terminología clínica y reduce el
	ruido, a la vez que produce vectores densos de baja dimensión que se almacenan e
	indexan eficientemente en una base vectorial. La implementación se apoya en la
	formulación de \cite{deerwester1990,manning2008} y en \texttt{scikit-learn}
	(\texttt{TruncatedSVD} sobre \texttt{TfidfVectorizer}).
	
	\subsection{Re-ordenamiento híbrido, BM25 y RAG}
	
	BM25 \cite{robertson2009} es una función de ranking probabilística que puntúa el
	solapamiento léxico exacto entre consulta y documento con saturación de
	frecuencia y normalización por longitud. Combinar una señal densa (LSI, orientada
	a \emph{recall} semántico) con una señal dispersa (BM25, orientada a la
	\emph{precisión} del término exacto) es la técnica estándar de recuperación
	híbrida. La Generación Aumentada por Recuperación (RAG) \cite{lewis2020}
	antepone una etapa de recuperación a un modelo generativo: el modelo redacta su
	respuesta condicionado a los documentos recuperados, lo que reduce las
	alucinaciones y permite citar fuentes. La retroalimentación de relevancia de
	Rocchio \cite{rocchio1971} re-pondera el vector de consulta acercándolo al
	centroide de los documentos marcados como relevantes y alejándolo de los no
	relevantes. Para diversificar recomendaciones se emplea \emph{Maximal Marginal
		Relevance} (MMR) \cite{carbonell1998}.
	
	\subsection{Evaluación}
	
	La calidad se mide con las medidas objetivas clásicas \cite{manning2008}:
	Precisión $P=|RR|/|RR\cup RI|$, Recobrado $R=|RR|/|RR\cup NR|$ y su media
	armónica $F_1=2PR/(P+R)$, donde $RR$, $RI$ y $NR$ son, respectivamente, los
	documentos relevantes recuperados, los irrelevantes recuperados y los relevantes
	no recuperados. Como $P$, $R$ y $F_1$ ignoran el orden, se añaden medidas
	sensibles al ranking: la ganancia acumulada descontada normalizada
	$\mathrm{NDCG}@k=\mathrm{DCG}@k/\mathrm{IDCG}@k$ con
	$\mathrm{DCG}@k=\sum_{i=1}^{k} \mathrm{rel}_i/\log_2(i+1)$
	\cite{jarvelin2002}, la precisión media (MAP) y el recíproco del rango medio
	(MRR).
	
	%==================================================================
	\section{Desarrollo e Implementación}\label{sec:desarrollo}
	%==================================================================
	
	\subsection{Arquitectura general}
	
	El sistema se implementa en Python como un \textbf{monolito modular} que
	combina los patrones \emph{Pipeline} y \emph{Microkernel}
	(Fig.~\ref{fig:arquitectura}). El flujo de una consulta es un \emph{pipeline} de
	etapas fijas y obligatorias; los módulos opcionales se conectan como
	\emph{plugins} a través de \emph{hooks} declarativos (\texttt{pre\_retrieval},
	\texttt{post\_retrieval}, \texttt{post\_ranking}) gestionados por
	\texttt{core/plugin\_pipeline.py}. Esta decisión arquitectónica ---frente a
	microservicios o \emph{Clean Architecture} pura, descartados por su sobrecarga
	para un equipo pequeño--- aporta tres beneficios: los plugins no contaminan el
	código obligatorio (si un plugin no se registra, el pipeline simplemente lo
	ignora), cada módulo implementa una interfaz abstracta (\texttt{core/interfaces.py})
	y es testeable de forma aislada, y todo corre en un único proceso, reproducible
	con Docker.
	
	\begin{figure}[t]
		\centering
		\begin{tikzpicture}[node distance=4.2mm and 6mm]
			\node[stage] (ui) {Interfaz Web\\(FastAPI + SPA)};
			\node[stage, below=of ui] (parse) {Parseo + Corrección\\ortográfica (Trie)};
			\node[hook, below=of parse] (h1) {pre\_retrieval\\(expansión de consulta)};
			\node[stage, below=of h1] (ret) {Recuperación LSI\\+ fallback web};
			\node[hook, below=of ret] (h2) {post\_retrieval};
			\node[stage, below=of h2] (rank) {Re-ranking híbrido\\(LSI + BM25)};
			\node[hook, below=of rank] (h3) {post\_ranking\\(recomendación)};
			\node[stage, below=of h3] (rag) {RAG\\(generación + evaluación)};
			
			\node[store, right=18mm of ret] (chroma) {ChromaDB\\(vectores LSI)};
			\node[store, below=4mm of chroma] (docstore) {Document\\Store};
			
			\draw[flow] (ui)--(parse);
			\draw[flow] (parse)--(h1);
			\draw[flow] (h1)--(ret);
			\draw[flow] (ret)--(h2);
			\draw[flow] (h2)--(rank);
			\draw[flow] (rank)--(h3);
			\draw[flow] (h3)--(rag);
			\draw[flow,<->] (ret)--(chroma);
			\draw[flow,<->] (ret.east)to[out=0,in=180](docstore.west);
			\draw[flow] (rag.east)to[out=0,in=-90](docstore.south);
			
			\begin{scope}[on background layer]
				\node[draw, dotted, rounded corners, fit=(parse)(rag), inner sep=4mm,
				label={[font=\scriptsize\itshape]right:Núcleo (pipeline + microkernel)}] {};
			\end{scope}
		\end{tikzpicture}
		\caption{Arquitectura Pipeline+Microkernel. Las etapas en azul son obligatorias;
			los recuadros verdes punteados son \emph{hooks} donde se enganchan plugins
			opcionales; los cilindros son los dos niveles de almacenamiento.}
		\label{fig:arquitectura}
	\end{figure}
	
	\subsection{Adquisición de datos: Crawler}
	
	La adquisición la realiza un \emph{crawler} genérico
	(\texttt{modules/crawler}) que separa estrictamente la lógica de E/S de la
	extracción específica de cada sitio. Las fuentes elegidas son tres servicios de
	salud de reputación contrastada: \textbf{MedlinePlus} (Biblioteca Nacional de
	Medicina de EE.\,UU., con contenido en español e inglés), \textbf{Mayo Clinic} y
	el \textbf{NHS} del Reino Unido. La confiabilidad y actualidad se aseguran por la
	naturaleza institucional y curada de estas fuentes y porque cada documento
	conserva en sus metadatos la fecha de última modificación declarada por el sitio.
	
	El descubrimiento de URLs es \emph{exclusivamente por sitemaps}: el crawler
	obtiene los \texttt{sitemap.xml} declarados por cada \emph{scraper}, expande
	recursivamente los \texttt{sitemapindex} (hasta una profundidad máxima de~3) y
	construye una cola de pares (URL, scraper). Las semillas, por tanto, son los
	sitemaps de los tres dominios. \emph{La búsqueda no sale del dominio del conjunto
		semilla}: al no implementarse seguimiento de enlaces y al filtrar cada URL con
	\texttt{can\_handle()} contra el dominio del scraper, el alcance queda acotado a
	las fuentes curadas, lo que es deseable para mantener la calidad del corpus. El
	respeto de las políticas de \emph{crawling} es explícito: se consulta y cachea el
	\texttt{robots.txt} de cada dominio (\texttt{RobotFileParser.can\_fetch}, con
	política \emph{fail-open}), se impone un \emph{rate limiting} por dominio
	(retardo mínimo configurable de 2~s entre peticiones), se aplica una política de
	reintentos con \emph{backoff} exponencial ante códigos 429/5xx, un \emph{timeout}
	de 15~s, un \texttt{User-Agent} identificable (\texttt{MedSRIBot/1.0 (proyecto
		académico UH)}) y un límite opcional \texttt{max\_pages}. Cada documento recibe un
	identificador estable \texttt{UUID5} derivado de su URL, lo que hace la
	adquisición idempotente entre re-crawls.
	
	\subsection{El corpus y su representatividad}
	
	El corpus indexado es \emph{textual}. Se compone de los documentos JSONL
	producidos por el crawler en \texttt{data/raw/} y de capítulos de libros médicos
	pre-procesados a JSON (ingeridos página a página, con identificadores del tipo
	\texttt{<libro>\_p<NNN>}). La suficiencia y representatividad se persiguen por construcción: combinar tres
	fuentes institucionales complementarias (enciclopedia médica, contenido para
	pacientes y guías clínicas) cubre tanto el registro divulgativo como el técnico,
	y el carácter bilingüe (es/en) amplía la cobertura terminológica. El recuento
	exacto de documentos depende de la ejecución del crawler; el sistema reporta en
	todo momento, mediante el endpoint \texttt{/api/stats}, el número de documentos
	originales, de \emph{chunks}, el tamaño del vocabulario y la dimensión latente
	$k$ efectivamente usada.
	
	\subsection{Indexación: índice invertido y normalización}
	
	La indexación (\texttt{modules/indexer}) construye un \textbf{índice invertido}
	que mapea cada término a una lista de \emph{postings} $(\text{idx\_doc},
	\text{tf})$. Antes, cada documento atraviesa el \texttt{TextProcessor}
	(\texttt{modules/text\_processor}), que aplica: normalización Unicode y a
	minúsculas, eliminación de caracteres no alfanuméricos preservando los propios
	del español (áéíóúüñ), tokenización y lematización con \texttt{spaCy}
	(\texttt{es\_core\_news\_md}), eliminación de \emph{stopwords} (NLTK español más
	una lista de \emph{stopwords} de dominio médico) y filtrado por longitud
	(2--20 caracteres). Los documentos largos se fragmentan previamente con un
	\texttt{TextChunker} (ventanas fijas de 300 palabras con solape de 50), de modo
	que cada \emph{chunk} es una unidad recuperable independiente que conserva en sus
	metadatos el \texttt{original\_doc\_id}. El vocabulario se ordena
	alfabéticamente (columnas deterministas) y se descartan los términos con
	frecuencia total inferior a~2, filtrando ruido y erratas. Como subproducto, los
	tokens lematizados de los documentos alimentan el vocabulario del corrector
	ortográfico (un \textbf{Trie} con distancia de Levenshtein) que más tarde corrige
	las consultas.
	
	\subsection{Modelo LSI, embeddings y base de datos vectorial}
	
	Sobre el índice invertido, \texttt{TfidfProcessor} construye la matriz TF--IDF y
	\texttt{LSIModel} (envoltorio de \texttt{TruncatedSVD}, $k=100$,
	\texttt{n\_iter}=7, semilla fija) la reduce al espacio latente. \textbf{Los
		vectores latentes resultantes son los \emph{embeddings} del sistema}: no se
	emplea un modelo de \emph{embeddings} neuronal externo, sino la propia proyección
	LSI, coherente con el modelo de recuperación elegido y reproducible sin
	dependencias de servicios en la nube. Estos vectores se almacenan en
	\textbf{ChromaDB}, configurado con espacio coseno (\texttt{hnsw:space=cosine}); la
	búsqueda por similitud usa el índice aproximado HNSW de ChromaDB, eficiente para
	hallar los vecinos más cercanos en alta dimensión sin recorrer toda la colección.
	
	El almacenamiento es de \textbf{dos niveles}: ChromaDB guarda únicamente IDs,
	vectores y metadatos mínimos (URL), mientras que el contenido completo reside en
	un \texttt{FileSystemDocumentStore}. Esta separación evita inflar la base
	vectorial, permite actualizar contenido sin re-indexar vectores y optimiza cada
	almacén para su propósito. Para información nueva, el sistema implementa
	\emph{indexación dinámica} (\emph{folding-in}, Fig.~\ref{fig:flujo}): un documento
	nuevo se vectoriza con la TF--IDF e IDF ya fijadas y se proyecta sobre la base
	latente existente sin reentrenar la SVD ($\mathbf{d}_{\text{proj}}=\mathbf{d}\,V_k$),
	y luego se anexa a ambos almacenes. Cuando la fracción de documentos añadidos por
	\emph{folding-in} supera el 20\,\% se dispara un re-ajuste completo
	(``balanceo'') que recaptura el vocabulario nuevo.
	
	\subsection{Flujo \emph{end-to-end} de una consulta}
	
	La Fig.~\ref{fig:flujo} narra el recorrido completo. (1)~El usuario escribe en la
	interfaz web y la petición llega a \texttt{POST /api/query}. (2)~\texttt{Indexer.
		build\_query} normaliza y \emph{corrige ortográficamente} la consulta contra el
	vocabulario del Trie; un mecanismo de interfaz permite además elegir entre la
	consulta corregida o la original tal cual se escribió. (3)~En el \emph{hook}
	\texttt{pre\_retrieval}, el plugin de expansión enriquece la consulta con
	sinónimos médicos. (4)~El \texttt{LSIRetriever} vectoriza con TF--IDF (descartando
	términos fuera de vocabulario), proyecta al espacio latente y realiza la
	\emph{recuperación en dos fases}: búsqueda vectorial en ChromaDB
	$\rightarrow$ IDs+scores, y recuperación del texto completo en el Document Store.
	(5)~Si existen juicios de relevancia previos para esa consulta, se aplica Rocchio
	sobre el vector latente. (6)~El \texttt{FallbackRetriever} decide si activar la
	búsqueda web. (7)~Tras \texttt{post\_retrieval}, el \texttt{HybridRanker}
	re-ordena por puntuación híbrida. (8)~En \texttt{post\_ranking} se preparan las
	recomendaciones. (9)~El \texttt{RAGService} genera la respuesta y el
	\texttt{RAGEvaluator} la puntúa. El resultado (documentos, respuesta, métricas y
	sugerencias) se serializa y se renderiza.
	
	\begin{figure}[t]
		\centering
		\begin{tikzpicture}[node distance=3.6mm and 5mm]
			\node[stage] (q) {Consulta del usuario};
			\node[stage, below=of q] (spell) {Corrección ortográfica\\(Trie + Levenshtein)};
			\node[stage, below=of spell] (exp) {Expansión por tesauro\\(pre\_retrieval)};
			\node[stage, below=of exp] (tfidf) {TF--IDF transform\\$\rightarrow$ proyección LSI};
			\node[stage, below=of tfidf] (search) {Búsqueda coseno\\en ChromaDB (top-k)};
			\node[stage, below=of search] (rocchio) {Rocchio\\(si hay feedback)};
			\node[stage, below=of rocchio] (fb) {¿Insuficiente?\\$\rightarrow$ fallback web};
			\node[stage, below=of fb] (hr) {HybridRanker\\$0.6\cdot$LSI$+0.4\cdot$BM25};
			\node[stage, below=of hr] (rag) {RAG (Groq) +\\evaluación de calidad};
			\node[stage, below=of rag] (out) {Render: tarjetas +\\respuesta + métricas};
			
			\foreach \a/\b in {q/spell, spell/exp, exp/tfidf, tfidf/search, search/rocchio,
				rocchio/fb, fb/hr, hr/rag, rag/out}
			\draw[flow] (\a)--(\b);
		\end{tikzpicture}
		\caption{Flujo \emph{end-to-end} de una consulta y módulos activados en orden.}
		\label{fig:flujo}
	\end{figure}
	
	\subsection{Detección de información insuficiente y búsqueda web}
	
	El \texttt{FallbackRetriever} decide que ``lo indexado no satisface la
	necesidad'' mediante tres criterios combinados: (i)~\emph{conteo}, cuando la
	recuperación local devuelve menos de \texttt{min\_results}=3 documentos;
	(ii)~\emph{puntuación}, cuando el mejor resultado local queda por debajo de
	\texttt{min\_score}=0.35; y (iii)~\emph{recencia}, cuando la consulta contiene
	marcas temporales (años 2020--2029, ``última hora'', ``nuevo tratamiento'',
	``actualizado'', etc.), pues el corpus curado puede no reflejar lo más reciente.
	Al activarse, el \texttt{InternetSearchRetriever} consulta DuckDuckGo (biblioteca
	\texttt{ddgs}), descarga y limpia las páginas con \texttt{WebContentFetcher}
	(\texttt{requests}+\texttt{BeautifulSoup}), asigna puntuaciones decrecientes por
	posición y marca los documentos con \texttt{origin="web"}. Los resultados web se
	\emph{fusionan} con los locales deduplicando por \texttt{doc\_id} (conservando la
	mayor puntuación) y se ordenan por score; además se \emph{cachean} en el
	Document Store para reutilizarse en consultas idénticas posteriores, evitando
	nuevas llamadas de red.
	
	\subsection{Re-ranking, factores de ordenamiento y posicionamiento visual}
	
	El \texttt{HybridRanker} produce el orden final fusionando dos señales
	normalizadas a $[0,1]$ por min--max: la similitud latente LSI (peso 0.6) y la
	puntuación léxica BM25 (peso 0.4), de modo que ningún criterio domine por escala.
	Además de la relevancia, intervienen otros factores: la \emph{frescura}
	(consultas marcadas como recientes desvían hacia resultados web actuales), el
	\emph{tipo/origen de contenido} (local frente a web, señalizado con insignias) y
	el \emph{perfil del usuario}, que el recomendador traduce en un \emph{boost}
	multiplicativo conservador (1.05--1.10) hacia las fuentes preferidas de cada
	perfil (p.\,ej.\ MedlinePlus para pacientes y cuidadores; Mayo Clinic/NHS para
	profesionales). Dos consultas equivalentes con perfiles distintos, o una consulta
	con marca de recencia frente a otra sin ella, producen así ordenamientos
	observablemente diferentes.
	
	En la interfaz (Fig.~\ref{fig:ui}), el posicionamiento responde a estas
	decisiones de diseño: los documentos se presentan como tarjetas ordenadas de
	arriba a abajo por puntuación híbrida descendente ---el primero es el más
	relevante---, cada tarjeta muestra título, fuente enlazable, un \emph{snippet}
	sensible a la consulta (la ventana de texto con mayor densidad de términos
	buscados) y un indicador porcentual de relevancia; una insignia distingue
	visualmente el contenido \emph{Local} del \emph{Web}. La respuesta generada por
	RAG ocupa un panel lateral propio para separar claramente ``documentos fuente''
	de ``respuesta sintetizada'', y un panel inferior de recomendaciones
	(``También te puede interesar'') agrupa contenido nuevo relacionado sin
	interferir con el ranking principal.
	
	
\begin{figure}[t]
	\centering
	\includegraphics[width=1\linewidth, height=.3\textheight]{"figuras/ShealthRI UI"}
	\caption{Interfaz de ShealtRI: barra de búsqueda con conmutadores (corrección ortográfica e inclusión de resultados web), tarjetas de resultados ordenadas por relevancia con insignias de origen, y panel lateral con la respuesta RAG.}
	\label{fig:shealthri-ui}
\end{figure}
	
	\subsection{Generación aumentada por recuperación (RAG)}
	
	El \texttt{RAGService} produce una respuesta adaptada a uno de seis perfiles de
	usuario (paciente, estudiante de medicina, profesional, asistente de diagnóstico,
	medicina natural y cuidador), cada uno con su propia plantilla de \emph{system
		prompt}, tono y áreas de enfoque. El proceso es: (1)~se re-ordenan los documentos
	con el \texttt{HybridRanker}; (2)~se construye un \emph{bloque de contexto} con a
	lo sumo los \textbf{3 documentos} mejor posicionados, recortando cada uno a
	800~caracteres ---es decir, \emph{no} se pasa todo lo recuperado, sino lo
	priorizado por relevancia---; (3)~se renderiza el \emph{prompt} insertando
	consulta y contexto; (4)~se invoca el modelo alojado en Groq
	(\texttt{llama-3.1-8b-instant}, \texttt{temperature}=0.3, \texttt{top\_p}=0.9,
	\texttt{max\_tokens}=1024). La \textbf{mitigación de alucinaciones} es
	estructural: el \emph{prompt} ancla explícitamente la respuesta a la sección
	``Información de referencia: \{context\}'' e instruye al modelo a \emph{indicar
		expresamente cuando los documentos de referencia son insuficientes}; la baja
	temperatura reduce la divagación. Si no hay clave de API o la llamada falla, el
	sistema degrada con elegancia a una \emph{respuesta de plantilla} que presenta
	directamente los \emph{snippets} recuperados. Por último, un \texttt{RAGEvaluator}
	puntúa cada respuesta en fidelidad, anclaje (\emph{groundedness}) y relevancia del
	contexto, métricas que la interfaz expone como indicador de calidad.
	
	\subsection{Plugins opcionales: expansión, retroalimentación y recomendación}
	
	Coherentes con el microkernel, los módulos opcionales se enganchan sin tocar el
	núcleo. La \textbf{expansión de consultas} (\texttt{plugins/expansion}, hook
	\texttt{pre\_retrieval}) busca cada término en un tesauro médico y añade sus
	sinónimos, abreviaturas y formas coloquiales al índice invertido de la consulta,
	filtrando los que no estén en el vocabulario del corpus y limitando a 6 términos
	para evitar la \emph{deriva de consulta}. La \textbf{retroalimentación de
		relevancia} (\texttt{plugins/feedback}) persiste los juicios del usuario
	(pulgar arriba/abajo en cada tarjeta) y, en la siguiente búsqueda de la misma
	consulta, aplica Rocchio (\(\alpha=1.0,\ \beta=0.75,\ \gamma=0.15\)) sobre el
	vector latente, acercándolo al centroide de lo relevante. El \textbf{recomendador}
	(\texttt{modules/recommender}) parte de los primeros resultados como semilla,
	calcula un centroide en el espacio LSI, recupera candidatos por coseno y los
	re-ordena con MMR ($\lambda=0.7$) para que la lista sea afín pero diversa,
	colapsando además páginas del mismo libro para no mostrar duplicados.
	
	%==================================================================
	\section{Evaluación y Resultados}\label{sec:evaluacion}
	%==================================================================
	
	El módulo de evaluación (\texttt{modules/evaluation}) implementa, como funciones
	puras y testeables, todas las medidas de la Sección~\ref{sec:teoria}: $P$, $R$,
	$F_1$ y $F_\beta$ general, $P@k$, $R@k$, \emph{Fallout}@k, MRR, MAP, DCG y
	NDCG@k. Sobre un conjunto de 10 consultas de prueba de dominio biomédico
	(\texttt{eval\_queries.json}) con sus juicios de relevancia
	(\texttt{eval\_qrels.json}), el endpoint \texttt{/api/eval} ejecuta el
	recuperador LSI y reporta tanto las métricas agregadas como el desglose por
	consulta, mostrándolos en la sección ``Evaluación'' de la interfaz
	(Tabla~\ref{tab:metricas}). El recobrado del corpus completo (\emph{Fallout}) se
	calcula contra el número de documentos \emph{originales} ---no de \emph{chunks}---
	para que el denominador sea coherente. Los valores numéricos concretos dependen
	del corpus efectivamente crawleado en cada despliegue; la infraestructura de
	evaluación permite reproducir la medición con un solo comando
	(\texttt{python cli.py --eval --eval-k 10}) y comparar LSI puro contra LSI con
	re-ranking híbrido (\texttt{--rerank}).
	
	\begin{table}[t]
		\centering
		\caption{Métricas implementadas en el módulo de evaluación y su rol.}
		\label{tab:metricas}
		\begin{tabular}{@{}lll@{}}
			\toprule
			\textbf{Métrica} & \textbf{Sensible al orden} & \textbf{Qué mide} \\
			\midrule
			Precisión / Recobrado / $F_1$ & No & Calidad del conjunto recuperado \\
			$P@k$ / $R@k$ / Fallout@k      & Parcial (corte $k$) & Calidad en el top-$k$ \\
			MAP                            & Sí & Precisión media sobre relevantes \\
			MRR                            & Sí & Posición del primer acierto \\
			NDCG@k                         & Sí & Ganancia graduada con descuento \\
			\bottomrule
		\end{tabular}
	\end{table}
	
	%==================================================================
	\section{Deficiencias, Conclusiones y Trabajo Futuro}\label{sec:conclusiones}
	%==================================================================
	
	ShealtRI demuestra que un SRI médico completo ---adquisición, indexación
	semántica, recuperación híbrida, fallback web y RAG--- puede construirse como un
	monolito modular reproducible, manteniendo cada responsabilidad aislada y
	testeable. No obstante, detectamos varias deficiencias y sus posibles soluciones.
	Primero, la base latente de LSI es \emph{estática}: los documentos añadidos por
	\emph{folding-in} no aportan vocabulario nuevo hasta el ``balanceo'', y los
	documentos web recuperados se cachean pero no se integran al espacio
	latente en la misma consulta; un trabajo futuro sería re-ajustar la SVD de forma
	incremental o adoptar \emph{embeddings} neuronales (\emph{sentence-transformers})
	que toleren vocabulario abierto. Segundo, el corpus puede quedar
	desactualizado frente a la evidencia médica reciente; integrar fuentes con
	\emph{feeds} fechados y un re-crawl programado lo mitigaría. Tercero, el Trie del
	corrector no se contrae al eliminar documentos, por lo que puede sugerir términos
	ya no indexados. Cuarto, el conjunto de evaluación es pequeño (10 consultas);
	ampliarlo con juicios de relevancia graduados por especialistas daría medidas más
	robustas. Finalmente, el sistema es solo textual: la recuperación multimodal de
	imágenes médicas queda como extensión.
	
	Bajo condiciones idílicas y empezando de cero, mantendríamos el patrón
	Pipeline+Microkernel ---que demostró su valor para entregas incrementales--- pero
	adoptaríamos desde el inicio \emph{embeddings} densos neuronales como
	representación primaria (con LSI como línea base comparativa), un índice vectorial
	con actualización verdaderamente incremental, y evaluación amplia y
	validada clínicamente. La separación de almacenamiento en dos niveles, la fusión
	híbrida y el anclaje estricto del RAG se conservarían sin cambios, por haber
	resultado decisiones acertadas.
	
	%==================================================================
	\begin{thebibliography}{10}
		
		\bibitem{manning2008} Manning, C.D., Raghavan, P., Schütze, H.: Introduction to
		Information Retrieval. Cambridge University Press, Cambridge (2008)
		
		\bibitem{baeza1999} Baeza-Yates, R., Ribeiro-Neto, B.: Modern Information
		Retrieval. ACM Press / Addison-Wesley, New York (1999)
		
		\bibitem{salton1983} Salton, G., McGill, M.J.: Introduction to Modern Information
		Retrieval. McGraw-Hill, New York (1983)
		
		\bibitem{deerwester1990} Deerwester, S., Dumais, S.T., Furnas, G.W., Landauer,
		T.K., Harshman, R.: Indexing by Latent Semantic Analysis. J. Am. Soc. Inf. Sci.
		41(6), 391--407 (1990)
		
		\bibitem{robertson2009} Robertson, S., Zaragoza, H.: The Probabilistic Relevance
		Framework: BM25 and Beyond. Found. Trends Inf. Retr. 3(4), 333--389 (2009)
		
		\bibitem{rocchio1971} Rocchio, J.J.: Relevance Feedback in Information Retrieval.
		In: Salton, G. (ed.) The SMART Retrieval System: Experiments in Automatic
		Document Processing, pp. 313--323. Prentice-Hall, Englewood Cliffs (1971)
		
		\bibitem{carbonell1998} Carbonell, J., Goldstein, J.: The Use of MMR,
		Diversity-Based Reranking for Reordering Documents and Producing Summaries. In:
		Proc. 21st ACM SIGIR, pp. 335--336. ACM, New York (1998)
		
		\bibitem{lewis2020} Lewis, P., et al.: Retrieval-Augmented Generation for
		Knowledge-Intensive NLP Tasks. In: Advances in Neural Information Processing
		Systems (NeurIPS), vol. 33, pp. 9459--9474 (2020)
		
		\bibitem{jarvelin2002} Järvelin, K., Kekäläinen, J.: Cumulated Gain-Based
		Evaluation of IR Techniques. ACM Trans. Inf. Syst. 20(4), 422--446 (2002)
		
		\bibitem{chroma} Chroma: the open-source embedding database,
		\url{https://www.trychroma.com}
		
		\bibitem{spacy} Honnibal, M., Montani, I.: spaCy: Industrial-Strength Natural
		Language Processing in Python, \url{https://spacy.io}
		
		\bibitem{sklearn} Pedregosa, F., et al.: Scikit-learn: Machine Learning in
		Python. J. Mach. Learn. Res. 12, 2825--2830 (2011)
		
		\bibitem{medlineplus} U.S. National Library of Medicine: MedlinePlus,
		\url{https://medlineplus.gov}
		
	\end{thebibliography}
	
	
\end{document}
